\section{Prefix to Postfix Conversion}
\label{sec:sq-prefix-to-postfix}

\problemstatement{Given a prefix expression, convert it to postfix
notation.}

\begin{patternbox}
The last of the six conversions, and by now a pure exercise in combining
two already-derived ideas: scan \textbf{right to left} (as in
Section~\ref{sec:sq-prefix-to-infix}, since the source is prefix), and
combine each pair as $\textit{left} + \textit{right} + \textit{op}$ (as in
Section~\ref{sec:sq-postfix-to-prefix}'s combination style, just operator
\emph{last} instead of first, and no parentheses needed).
\end{patternbox}

\subsection*{Understanding the problem carefully}

No new idea is required: take the right-to-left scan direction from
Section~\ref{sec:sq-prefix-to-infix} (because the source is prefix) and
the parenthesis-free, operator-position-only combination style appropriate
for postfix output (operator written last).

\begin{ideabox}[Combine as Postfix Instead of Infix]
Scan right to left. Push every operand. On operator $\textit{op}$: pop
$\textit{left}$, pop $\textit{right}$ (same pop order as
Section~\ref{sec:sq-prefix-to-infix}, since the source notation and scan
direction are unchanged); push $\textit{left} + \textit{right} +
\textit{op}$.
\end{ideabox}

\subsection*{Algorithm}

\begin{enumerate}
  \item Initialize an empty stack of strings.
  \item For each token, scanning right to left: operand $\to$ push.
        Operator $\textit{op}$ $\to$ pop $\textit{left}$, pop
        $\textit{right}$; push $\textit{left} + \textit{right} +
        \textit{op}$.
  \item The stack's single remaining element is the postfix result.
\end{enumerate}

\subsection*{Dry run}

\dryrun{Take $\texttt{"*+abc"}$.}

Scanning right to left: \texttt{c}, \texttt{b}, \texttt{a}, \texttt{+},
\texttt{*}.

\begin{itemize}
  \item \texttt{c}: push. Stack $=[\texttt{c}]$.
  \item \texttt{b}: push. Stack $=[\texttt{c},\texttt{b}]$.
  \item \texttt{a}: push. Stack $=[\texttt{c},\texttt{b},\texttt{a}]$.
  \item \texttt{+}: pop left$=\texttt{a}$, pop right$=\texttt{b}$. Push
        \texttt{ab+}. Stack $=[\texttt{c},\texttt{ab+}]$.
  \item \texttt{*}: pop left$=\texttt{ab+}$, pop right$=\texttt{c}$.
        Push \texttt{ab+c*}.
\end{itemize}

Result: \texttt{ab+c*} -- exactly the postfix expression we started from
back in Section~\ref{sec:sq-postfix-to-infix}, confirming the whole
six-conversion group is internally consistent.

\subsection*{C++ implementation}

\begin{lstlisting}
#include <bits/stdc++.h>
using namespace std;

string prefixToPostfix(string s) {
    stack<string> st;

    for (int i = s.size() - 1; i >= 0; i--) {
        char c = s[i];
        if (isalnum(c)) {
            st.push(string(1, c));
        } else {
            string left = st.top(); st.pop();
            string right = st.top(); st.pop();
            st.push(left + right + string(1, c));
        }
    }
    return st.top();
}

int main() {
    cout << prefixToPostfix("*+abc") << endl;
    return 0;
}
\end{lstlisting}

\begin{complexitybox}
\textbf{Time Complexity.} $O(n^2)$ worst case with naive string
concatenation, $O(n)$ with more careful string building.
\textbf{Space Complexity.} $O(n)$.
\end{complexitybox}

\begin{takeawaybox}
Across all six conversions in this group, there are really only two
independent decisions: \emph{which direction to scan} (left to right if
the source is postfix or plain infix; right to left if the source is
prefix) and \emph{how to format the combination step} (parenthesized
infix, operator-first prefix, or operator-last postfix). Every one of the
six problems is a different combination of these same two choices.
\end{takeawaybox}
